% Generated draft: requires mathematical and source review.
\documentclass[11pt]{article}
\usepackage[margin=25mm]{geometry}
\usepackage{fontspec}
\setmainfont{Latin Modern Roman}
\usepackage{amsmath,amssymb,mathtools,unicode-math}
\setmathfont{Latin Modern Math}
\usepackage{enumitem,needspace,xurl,hyperref}
\setlength{\parindent}{0pt}
\setlength{\parskip}{5pt}
\setlength{\emergencystretch}{3em}
\newcommand{\N}{\mathbb{N}}
\newcommand{\Z}{\mathbb{Z}}
\newcommand{\Q}{\mathbb{Q}}
\newcommand{\R}{\mathbb{R}}
\newcommand{\C}{\mathbb{C}}
\newcommand{\F}{\mathbb{F}}
\newcommand{\A}{\mathbb{A}}
\newcommand{\PP}{\mathbb{P}}
\newcommand{\E}{\mathbb{E}}
\newcommand{\eps}{\varepsilon}
\newcommand{\dd}{\,\mathrm{d}}
\begin{document}
\section*{UnsolvedMath ID: 701350}
\textbf{Source problem code:} catalogue entry
\subsection{Definitions and mathematical statement}
Work in the first-order language of groups \(L_{\mathrm{grp}}=\{1,\cdot,{}^{-1}\}\). For a group \(H\), let \(\operatorname{Th}(H)\) denote its complete first-order theory in this language. A complete theory \(T\) is called stable if there is no formula \(\varphi(\bar x;\bar y)\), with fixed finite tuples of variables, having the order property: for every integer \(n\geq 1\), there are a model \(M_n\models T\) and tuples \(\bar a_1,\ldots,\bar a_n\), \(\bar b_1,\ldots,\bar b_n\) in \(M_n\) such that
\[
M_n\models\varphi(\bar a_i;\bar b_j)\quad\Longleftrightarrow\quad i<j
\qquad(1\leq i,j\leq n).
\]
For groups \(A\) and \(B\), their free product \(A*B\) is the group with canonical embeddings of \(A\) and \(B\) characterized by the following universal property: for every group \(K\) and homomorphisms \(f:A\to K\), \(g:B\to K\), there is a unique homomorphism \(A*B\to K\) extending both \(f\) and \(g\).

The question attributed to E. Jaligot is whether, for every pair of groups \(A,B\),
\[
\bigl(\operatorname{Th}(A)\text{ is stable}\ \land\ \operatorname{Th}(B)\text{ is stable}\bigr)
\quad\Longrightarrow\quad
\operatorname{Th}(A*B)\text{ is stable}.
\]
The supplied statement imposes no restrictions on cardinality or nontriviality. Its registry status is \emph{open}; this is a supplied classification, not an independently verified assessment of the literature.
\subsection{Short English statement}
Must the free product of any two groups whose complete first-order theories are stable also have a stable complete first-order theory? The problem is attributed to E. Jaligot (2010).
\subsection{Sources}
\begin{itemize}
\item Supplied local UnsolvedMath registry record: permanent ID \(701350\), source code \texttt{catalogue entry}, title \emph{Kourovka Notebook Problem 17.48}. The record attributes the problem to E. Jaligot, gives the proposed year as 2010, and labels it open.
\item \emph{The Kourovka Notebook (Unsolved Problems in Group Theory)}, original problem 17.48, 17th issue (2010). The supplied background also identifies the 21st issue edition (2026), \texttt{arXiv:1401.0300v48}. The supplied collection description names E. I. Khukhro and V. D. Mazurov as editors. These bibliographic details are supplied source attribution and have not been independently checked.
\end{itemize}
\end{document}
